% 本文件由 LaTeX 排版 Agent 根据有效 translation.json 自动生成。
% author=Augustine of Hippo; work_id=1307; title=向慕道者讲论信经

\chapter{\WorkTitle{向慕道者讲论信经}}

% structure: p0001 source=h1 target=omitted reason=page-title-already-rendered-by-parent

1. 孩子们，请接受信德准则，即称为\OriginalTerm{信经}{Symbol}（或信经）的准则。你们既已接受，就当写在心中，每日向自己诵念；在睡眠前，在出门前，以信经武装自己。没有人将信经写成可供阅读的形式；但为复习之故，免得遗忘消磨了关怀所传授的，让你们的记忆成为记录卷轴：你们将要听见的，就是你们要相信的；你们所相信的，将要用舌头诵出。因为宗徒说：\BibleQuote{心里相信，可成为义；口里承认，可得到救恩。}\newline{}这是你们将要诵念并回答的信经。你们所听见的这些话，分散在圣经各处，但从中收集并归纳为一，使慢记忆的人不致困难；使每个人都能说出、能持守他所信的。难道你们仅仅听见天主是全能的吗？但当你们由教会——你们的母亲——所生时，你们就开始以祂为父。\par

2. 因此，对于你们现在所领受、所默想、并默想后所持守的，你们应当说：\Quote{我信全能的天主父。}天主是全能的，然而，尽管全能，祂不能死，不能被欺骗，不能说谎；并且，如宗徒所说：\BibleQuote{祂不能否认自己。}有许多事祂不能做，却仍然是全能的！正因祂不能做这些事，所以才是全能的。因为如果祂能死，祂就不是全能的；如果祂能说谎、被欺骗、行不义，祂就不是全能的：因为如果这些在祂内，祂就不配称为全能。对我们的全能之父来说，犯罪是完全不可能的。祂行祂所愿的一切：那就是全能。祂行祂所愿的一切正当事、公义事；但凡是邪恶的事，祂不愿行。无人能抵抗全能者，使祂不施行祂的意愿。是祂创造了天、地、海及其中的万物，有形的和无形的。无形的如在天上的：宝座、宰制者、率领者、掌权者、总领天使、天使；这一切，若我们生活正直，将是我们共同的天乡公民。祂创造了天上可见之物：太阳、月亮、星辰。祂以地上的动物装饰大地，以飞鸟充满天空，以行走和爬行的充满陆地，以游鱼充满海洋：祂以各自适宜的受造物充满了一切。祂也按照自己的肖像和模样造了人，在于理智：因为天主的肖像就在理智内。这就是为什么理智甚至不能为自身所理解，因为它内有天主的肖像。我们被造是为管理其他受造物；但因原祖的罪，我们都堕落了，进入了死亡的继承。我们变得卑微，成为可朽的，充满恐惧和错误：这是罪所应得的；每个生来的人带着这应得和罪债。这就是为什么，如你们今天所见、所知，连幼童也要经历嘘驱魔和驱魔礼：为驱走他们内那欺骗人类以占有全人类的仇敌魔鬼的权势。所以，受驱魔或嘘驱魔的并非天主的受造物在婴儿身上，而是那掌权者，因为所有带着罪而生的人都在他权下；他是罪人之首。因这缘故，由于一人堕落并使众人陷入死亡，就派遣了一位无罪的，祂要把所有相信祂的人从罪中救出，带入生命。\par

3. 因此我们也相信祂的圣子，即全能天主父的\Quote{独生子，我们的主。}当你们听到天主的独生子时，要承认祂是天主。因为天主的独生子不可能是非天主。祂所是的，同样生了祂，尽管祂并不是祂所生的那位。如果祂真是圣子，祂就是父所是的；如果祂不是父所是的，祂就不是真圣子。看看有死和属地的受造物：每个生出与自身同类的。人不生牛，羊不生狗，狗不生羊。无论生子者是什么，它就生出什么。所以你们要勇敢、坚定、忠信地持守：天主父所生的就是祂自己，全能者。这些有死的受造物是通过腐朽而生的。天主是这样生的吗？有死的生于者生出与自身同类的；不朽者生出祂所是的：可朽的生出可朽的，不朽的生出不朽的：可朽的以可朽的方式，不朽的以不朽的方式；是的，祂这样生出祂所是的，一位生出一位，因此是独生子。你们知道，当我把信经念给你们时，我如此说了，你们也有义务相信：我们\Quote{信全能的天主父，并信耶稣基督，祂的独生子。}这里，当你们相信祂是独生子时，也要相信祂是全能的：因为不可认为天主父行祂所愿的，而天主子不行祂所愿的。父与子同一意志，因为同一性体。因为子的意志绝不可能与父的意志有任何分离。天主与天主，两者同是一位天主；全能者与全能者，两者同是一位全能者。\par

4. 我们不引进两位天主，如同有些人所说的：\Quote{天主父和天主圣子，但父是天主更大，子是天主较小}。他们二者是什么？两位天主么？你羞于说出口，也当羞于相信。你说：主天主父和主天主圣子；而圣子亲自说：\Quote{没有人能事奉两位主}（\BibleRef{玛 6:24}）。在祂的家中，我们难道要这样，如同在一个大宅里，有家主和他的儿子，我们就说：较大的主，较小的主？要远离这种想法。如果你在心中如此设想，你就是在\Quote{同一灵魂}里立偶像。彻底摒弃它。先相信，然后理解。如今，天主赐给谁，使他在相信后很快理解，那是天主的恩赐，而非人的软弱。尽管如此，如果你还不理解，要相信：一位天主父，天主基督是天主圣子。二者是什么？一位天主。那么二者如何被说成是一位天主？如何？你感到惊奇吗？在\BibleBook{宗徒}{大事录}中，记载说：\Quote{众信徒都一心一意}（\BibleRef{宗 4:32}）。有许多灵魂，信德使它们成为一体。那里有成千上万的灵魂；他们彼此相爱，许多人成为一体：他们在爱德之火中爱天主，从众多达到了美的合一。如果那么多灵魂因爱之亲密成为一体，那么在天主那里，没有差异，只有完全的平等，爱之亲密该是何等！如果在地上、在人间，能有如此大的爱，使这么多灵魂成为一个灵魂，那么圣父与圣子从未分离，他们岂能不是一位天主？只是，那些灵魂既可称为众多灵魂，也可称为一个灵魂；但天主，祂有不可言喻的最高结合，可称为一位天主，而不是两位天主。\par

5. 圣父做祂愿意做的事，圣子也做祂愿意做的事。不要想象一位全能的天主父和一位非全能的天主子：那是错误，要在你内涂抹它，不要让它附着在你的记忆中，不要让它被吸收入你的信德，如果你们中有人偶然把它吸进去了，就让他呕吐出来。圣父是全能的，圣子也是全能的。如果全能者没有生出全能者，祂就没有生出真实的圣子。因为，弟兄们，我们说什么呢？如果圣父更大，生了一个比祂小的圣子？我说什么？生育？人生育时更大，生的儿子更小：这是真的；但那是因为前者衰老，后者长大，并通过成长达到父亲的样式。天主子，如果祂不长大（因为天主也不能衰老），祂被生出时就是完美的。既然是完美的被生，如果祂不长大，且没有保持更小，祂就是平等的。为使你们知道全能者出自全能者，请听那位是真理者。真理本身说的话是真实的。真理说什么？圣子是真理，祂说什么？\Quote{凡父所做的事，子也照样做}（\BibleRef{若 5:19}）。圣子是全能的，做一切祂愿意做的事。因为如果圣父做一些圣子不做的事，圣子就虚假地说：\Quote{凡父所做的事，子也照样做}。但因为圣子说了真话，要相信：\Quote{凡父所做的事，子也照样做}，而你已经相信圣子是全能的。这句话虽然你在\WorkTitle{信经}中没有说出，但当你相信独生子、祂自己就是天主时，你表达的就是这个。圣父有什么是圣子没有的吗？这是亚流异端亵渎者说的，不是我。但我说什么？如果圣父有什么是圣子没有的，那么圣子在说\Quote{凡父所有的，都是我的}（\BibleRef{若 16:15}）时就是说谎了。许多无数的见证都证明，圣子是圣父天主的真实儿子，圣父天主有其独生子天主，圣父和圣子是一位天主。\par

6. 但这位全能父天主的独生子，让我们看看祂为我们做了什么，祂为我们受了什么。\Quote{因圣神由童贞玛利亚诞生}。祂如此伟大，与圣父平等，因圣神由童贞玛利亚诞生，诞生得卑微，为借此医治骄傲的人。人自高而堕落；天主自卑而提升他。基督的卑微是什么？天主向堕落的人伸出手。我们跌倒，祂下降；我们躺卧，祂俯身。让我们紧握并起来，以免落入惩罚。因此，祂向我们的俯身就是：\Quote{因圣神由童贞玛利亚诞生}。祂作为人的诞生本身，既是卑微的，也是崇高的。为何卑微？因为祂作为人由人所生。为何崇高？因为祂由童贞女所生。童贞女怀孕，童贞女生产，生产后仍是童贞女。\par

7. 接下来呢？\Quote{在般雀·比拉多执政时受难}。当时他任总督并是审判官，就是这位般雀·比拉多，在基督受难的时候。审判官的名字标明了时代，当祂在般雀·比拉多手下受难时：当祂受难时，\Quote{被钉在十字架上，死，而被埋葬}。谁？什么？为谁？谁？天主的独生子，我们的主。什么？被钉在十字架上，死，被埋葬。为谁？为不虔敬的人和罪人。极大的俯就，极大的恩宠！\Quote{我拿什么报答上主赐给我的一切恩惠？}（\BibleRef{咏 116:12}）\par

8. 祂在万世之前、在诸世界之前受生。\Quote{受生之前。}在什么之前？在祂内没有之前的那位？丝毫不要想象在基督从圣父受生的那个诞生之前有任何时间；我所说的是那个诞生，因着它祂是天主全能者的独生子、我们的主；我首先说的就是那个。不要在这个诞生中设想时间的开始；不要设想任何一段永恒，在其中圣父存在而圣子不存在。因为当圣父存在时，圣子就存在。而那个\Quote{当}是什么？那里没有开始？因此，圣父永远无始，圣子永远无始。你将会说：如果祂没有开始，祂是如何受生的？由永恒者，\OriginalTerm{同永恒者}{coeternal}。在任何时间里，圣父都不是存在而圣子不存在，然而圣子是由圣父所生。从哪里能找到任何相似物呢？我们处于地上的事物中，处于可见的受造物中。让大地给我一个相似物：它给不出。让水给我一个相似物：它没有可以用来给的。让某个动物给我一个相似物：它也不能做到。动物确实会生育，生育者和被生者都存在；但父亲在先，儿子在后出生。让我们找到\OriginalTerm{同岁者}{coeval}并设想它是同永恒的。如果我们能找到与儿子同岁的父亲，以及与父亲同岁的儿子，就让我们相信天主圣父与祂的圣子同岁，而天主圣子与祂的圣父同永恒。在地上，我们能找到一些同岁者，但找不到任何同永恒者。让我们延伸同岁者，并设想它为同永恒者。有人可能会让你紧张，说：\Quote{何时能找到与儿子同岁的父亲，或与父亲同岁的儿子？为了生育，父亲在年龄上先于；为了被生，儿子在年龄上后于：但这样的父亲与儿子同岁，或儿子与父亲同岁，怎么可能呢？}请你们想象火为父亲，它的光芒为儿子；看，我们找到了同岁者。从火开始存在的那个瞬间，那个瞬间它就生出光芒：不是火先于光芒，也不是光芒后于火。如果我们问：哪个生哪个？火生光芒，还是光芒生火？你们立刻凭自然感觉、凭你们内心天生的智慧都喊出：火生光芒，不是光芒生火。看，这里你有一个开始的父亲；看，一个同时开始的儿子，既不在前也不在后。看，那么这是一个开始的父亲，看，一个同时开始的儿子。如果我向你们展示了一个开始的父亲和一个同时开始的儿子，那么请相信圣父无始，并且与祂一起的圣子也无始；一位是永恒的，另一位是同永恒的。如果你们在学习上进步，你们就会明白：要努力进步。你们已经有了出生；但你们还应该有成长；因为没有人一开始就是完美的。至于天主之子，祂确实能够出生而完美，因为祂是在时间之外受生，与圣父同永恒，远在万有之前，不是在年龄上，而是在永恒中。祂然后受生为同永恒者，关于这个出生，先知说：\BibleQuote{他的世代谁能述说？}祂在时间之外由圣父受生，在时期圆满时由童贞玛利亚出生。这个诞生有之前的时间为其预备。在适当的时候，当祂愿意、当祂知道时，祂就诞生了：因为祂不是违背自己意愿诞生的。我们中没有一个是因自己的意愿而出生，也没有一个在自己愿意的时候死去：祂，当祂愿意时，就诞生了；当祂愿意时，就死去了：祂如自己所愿意的方式，由童贞玛利亚诞生；如自己所愿意的方式，死在十字架上。凡祂愿意的，祂都做了：因为祂是如此的人，同时无形的，祂是天主；天主摄取，人被摄取；同一位基督，天主和人。\par

9. 关于祂的十字架，我该说什么，讲什么呢？祂选择了这种最极端的死亡，以便祂的殉道者不致因任何种类的死亡而恐惧。祂在作为人的生活里显示出教义，在祂的十字架上彰显了忍耐的榜样。看，这里有工作：祂被钉十字架；工作的榜样：十字架；工作的奖赏：复活。祂在十字架上向我们展示了我们应忍受什么，在复活中展示了我们应该盼望什么。如同竞技比赛中的一位完美教练，祂说：\Quote{做，并承担；做工作并接受奖赏；在比赛中奋斗，你将被加冕。}工作是什么？服从。奖赏是什么？没有死亡的复活。我为什么加了\Quote{没有死亡？}因为拉匝禄复活了，却又死了；基督复活了，\BibleQuote{不再死亡，死亡不再统治祂。}\par

10. 圣经记载说：\BibleQuote{你们听见了约伯的忍耐，也看见了主的结局。}当我们读到约伯所受的巨大考验时，令人战栗，令人退缩，令人发抖。他得到了什么？他得到了所失之物的双倍。因此，人不可怀着对现世赏报的期望而愿意忍耐，并对自己说：\Quote{让我忍受损失，天主会还给我双倍的儿子；约伯得到了全部的双倍，并且生了和他埋葬的儿子一样多的儿子。}难道这不是双倍吗？是的，正是双倍，因为先前的儿子仍然活着。谁也不要说：\Quote{让我忍受灾祸，天主会像对待约伯那样报答我。}这样，忍耐就不再是忍耐，而是贪婪了。因为如果那位圣人没有忍耐，也没有勇敢地承受临到他身上的一切，那么主所给的见证从何而来？主说：\BibleQuote{你注意到我的仆人约伯没有？因为在地面上没有人像他那样，一个无可指摘的人，真正敬拜天主的人。}我的弟兄们，这位圣人在主面前得到了何等样的见证啊！然而，一个坏女人却试图用劝说欺骗他，她同样代表了那蛇，就像在乐园里蛇欺骗了天主最初创造的人一样，同样这里也通过建议亵渎，以为能欺骗一个取悦天主的人。他遭受了什么事啊，我的弟兄们！谁能遭受如此多的痛苦，在他的财产、房屋、儿子、身体上，甚至在他的妻子身上，她被留下来作他的试探者！至于留下的那个妻子，魔鬼早就想夺走她，但他留住她作为帮手：因为藉着厄娃他征服了第一个人，因此他留住了一个厄娃。那么，他遭受了什么事啊！他失去了所有的一切；他的房屋倒塌了；但愿仅此而已！它也压死了他的儿子们。而且，为了看到忍耐在他身上有多大的位置，请听他怎样回答：\BibleQuote{上主赐的，上主收回；随上主的意愿成就；愿上主的名受赞颂。}他取走了他给的，难道赐予者丢失了吗？他取走了他给的。好像他在说：他拿走了一切，让他拿走一切，打发我赤裸而去，但让我保有他。如果我有天主，我还会缺少什么？如果我没有天主，其他一切对我又有什么好处？然后轮到他的身体，他从头到脚受了伤；他成了一个溃烂的伤口，一堆爬行的蠕虫；然而在天主面前他表现出坚定不移，屹立不动。那女人，魔鬼的帮凶而非丈夫的安慰者，想怂恿他亵渎天主。她说：\Quote{你还要忍受到几时？你说句反对主的话，然后死掉吧。}所以，因为他被贬抑，他将被高举。主这样做是为了向人显示；至于他的仆人，他在天上为他预备了更大的事。所以，被贬抑的约伯，主高举了他；而被高举的魔鬼，主贬抑了他：因为\BibleQuote{他贬抑一个，高举另一个。}但是，我亲爱的弟兄们，当任何人遭受类似的磨难时，不要寻求现世的赏报：例如，如果他遭受损失，他或许不能说：\BibleQuote{上主赐的，上主收回；随上主的意愿成就；愿上主的名受赞颂。}却只想着得回双倍。让忍耐赞美天主，而不是贪婪。如果你失去的是为了得到双倍，并为此赞美天主，那么你是出于贪心赞美，而不是出于爱。不要以为这是那位圣人的榜样；你在欺骗自己。当约伯忍受一切时，他并不希望得到双倍。在他第一次失物时忍受损失并埋葬儿子尸体的宣告中，以及第二次他因身体疮痍而痛苦时，你可以看到我所说的。他第一次宣告的话这样说：\BibleQuote{上主赐的，上主收回；随上主的意愿成就；愿上主的名受赞颂。}他本可以说：\Quote{上主赐的，上主收回；那位取走的能再次赐予；能带回比他所取更多的。}但他没有这样说，而是说：\BibleQuote{随上主的意愿成就。}因为主乐意，就让我也乐意；不要让顺服的仆人不喜欢好主人所喜欢的；不要令病人不喜欢医生所喜欢的。听听他另一次宣告：他对妻子说：\BibleQuote{你说话像糊涂女人！我们若从天主那里接受恩惠，难道不接受灾祸吗？}他没有加上——尽管如果他加上，也是真的——\BibleQuote{主既能将我的身体恢复原状，也能将从我们这里取走的加多几倍。}免得他被人视为因盼望这个而忍耐。他不是这样说，也不是这样盼望。但是，为了教训我们，主为他做了这事，虽然他并不盼望，由此我们被教训：天主与他同在；因为如果主没有恢复那些东西，固然有冠冕隐藏着，我们看不见。因此，神圣的圣经在劝勉忍耐和对未来事物的盼望——而非对现世赏报的盼望——时说什么？\BibleQuote{你们听见了约伯的忍耐，也看见了主的结局。}为什么是\Quote{约伯的忍耐}，而不是\Quote{你们看见了约伯自己的结局}？你会张开口等着\Quote{双倍}，会说：\BibleQuote{感谢天主，让我忍受吧：我会得到双倍，像约伯那样。} \OriginalTerm{约伯的忍耐}{the patience of Job}，\OriginalTerm{主的结局}{the end of the Lord}。我们认识约伯的忍耐，也认识主的结局。主的什么结局？\BibleQuote{我的天主，我的天主，你为什么舍弃了我？}这是主悬在十字架上说的话。他似乎为了现世的福乐而舍弃了他，却没有舍弃他永恒的不朽。这就是\Quote{主的结局}。犹太人抓住他，嘲笑他，捆绑他，给他戴上荆棘冠冕，用唾沫侮辱他，鞭打他，辱骂他，把他挂在木架上，用枪刺他，最后埋葬了他。他仿佛是舍弃了。但被谁舍弃？被那些侮辱他的人。因此，你也要忍耐到底，以便复活而不死，即永远不死，如同基督一样。因为正如我们读到：\BibleQuote{因为基督从死者中复活，就不再死。}\par

11. \Quote{祂升了天堂}：你要相信。\Quote{祂坐在圣父之右}：你要相信。\Quote{坐}应理解为居住，如同我们说某个人，\Quote{他在那个地方住了（\OriginalTerm{住了}{sedit}）三年}。圣经也有这样的表述，说某人在某城\OriginalTerm{居住过}{sedisse}一段时间。难道这意味着他坐下后就从未起来过吗？因此，人的居所被称为\OriginalTerm{座位}{sedes}。人们\Quote{坐}在这里时，难道他们总是一直坐着吗？难道没有站起、行走或躺下吗？然而它们仍被称为座位。如此，你当相信基督在圣父之右的居住：祂就在那里。不要让你的心问你说：\Quote{祂在做什么？}不要寻求那不允许你找到的东西：祂在那里，这对你已足够。祂是蒙福的，而那称为圣父之右的福乐——这至福本身的名字就是圣父之右——因为如果我们按照肉身来理解，那么既然祂坐在圣父之右，圣父将在祂的左边。若把圣子在右、圣父在左这样放在一起，难道合乎虔敬吗？那里全是右边，因为那里没有痛苦。\par

12. \Quote{从那里祂要来审判生者死者。}生者，即那些活着并存留的人；死者，即那些已经先行离去的人。也可以这样理解：生者即义人，死者即不义之人。因为祂要审判二者，各按其报。在审判中，祂将对义人说：\BibleQuote{我父所祝福的，你们来吧！承受自创世以来为你们预备的国度。}为此你们预备自己，为此你们盼望，为此你们生活，并如此生活，为此你们相信，为此你们领受圣洗，以便可以对你们说：\BibleQuote{我父所祝福的，你们来吧！承受自创世以来为你们预备的国度。}对左边的人，祂说什么？\BibleQuote{进入那为魔鬼和他的使者预备的永火里去。}如此，基督将审判生者死者。我们已经论及基督的第一次诞生，那是超越时间的；论及另一次诞生，在时期圆满之时，基督由童贞女诞生；论及基督的受难；论及基督的审判来临。关于基督、天主的独生子、我们的主，应当论述的都已经论述了。但圣三尚未完全。\par

13. 信经接着说：\Quote{也信圣神。}这圣三，一位天主，一个性体，一个本体，一个德能；最高平等，无分割，无差异，永恒的爱之亲密。你们愿意知道圣神是天主吗？领受圣洗，你们将成为祂的殿宇。宗徒说：\BibleQuote{你们不知道你们的身体是在你们内由天主所得的圣神的殿吗？}殿宇是为天主的；同样，所罗门君王兼先知，奉命为天主建造殿宇。若他建造殿宇给太阳、月亮、星辰或某位天使，天主岂不要定他的罪吗？因此，因为他建造殿宇给天主，表明他敬拜天主。他用什么建造呢？用木石，因为天主愿意通过祂的仆人在世上为自己造一所房屋，在那里可以被祈求，被纪念。对此，荣福斯德望说：\BibleQuote{所罗门为祂建造了殿宇，但至高者不住在人手所造的殿宇中。}那么，如果我们的身体是圣神的殿宇，那为圣神建造殿宇的是怎样的天主呢？但那是天主。因为若我们的身体是圣神的殿宇，那么为圣神建造这殿宇的那位，就是建造我们身体的那位。请听宗徒说：\BibleQuote{天主将此身体调和，将更大的光荣加给那有缺欠的肢体}，当时他正在谈论不同的肢体，使身体中没有分裂。天主创造了我们的身体。青草，是天主创造的；我们的身体是谁创造的？我们如何证明青草是天主的创造？那给青草披上衣服的，就是创造它的。请读福音：\BibleQuote{田野的草，今天还在，明天就被投入炉中，天主尚且这样装扮它}，它说。那么，创造它的那位就是装扮它的那位。宗徒又说：\BibleQuote{糊涂人哪！你所种的，若不先死，决不会生；并且你所种的，不是那将来的形体，而是赤裸的种子，好比麦粒或别的种子；但天主随自己的心意给它一个形体，并给每种种子以各自的形体。}既然如此，若是天主建造我们的身体，天主建造我们的肢体，而我们的身体是圣神的殿宇，那么你们不要怀疑圣神是天主。也不要把祂当作第三位神；因为父、子和圣神是一位天主。你们应当这样相信。\par

14. 在赞美圣三之后，接着说：\Quote{圣教会。}指出天主和祂的殿宇。\BibleQuote{因为天主的殿宇是圣的，}宗徒说，\BibleQuote{你们就是这殿宇。}这同一教会就是圣教会，唯一的教会，真教会，公教会，与一切异端争战：它能争战，但不能被战胜。至于异端，它们都从教会中出去，如同无用的枝条从葡萄树上剪下；但教会本身留在自己的根上，留在自己的葡萄树上，留在自己的爱德上。\BibleQuote{地狱的门不能胜过她。}\par

15. \Quote{罪过的赦免。}你领受圣洗时，信经的[这一条]就完全在你内。谁也不可说：\Quote{我犯了这样那样的罪，恐怕那罪不得赦免。}你做了什么？你犯了多大的罪？说出你行过的任何可憎之事，沉重、可怕，甚至一想就战栗：随你做了什么都行：你杀了基督吗？没有比那更恶劣的行为，因为也没有比基督更好的。杀基督是多么可怕的事！然而犹太人杀了祂，后来许多人信了祂，喝了祂的血：他们所做的罪已得赦免。当你受洗后，要在天主的诫命中持守善生，好能守护你的圣洗直到终末。我不是说你在此世不会犯罪；但那些是小罪，没有它们此生便不存在。圣洗是为了一切罪而设立的；为了我们不可避免的轻罪，祈祷被设立。祈祷说了什么？\Quote{宽免我们的罪债，如同我们也宽免欠我们债的人。}我们在圣洗中一次得洁净，我们在祈祷中每天得洁净。只是，不要犯那些必须与基督身体分离的事——愿这事远离你！因为你所见做补赎的人，是行了可憎之事，或是奸淫，或是某种大恶：为此他们做补赎。因为如果他们的罪是轻的，每日的祈祷就足以涂去。\par

16. 可见，在教会内罪过得以赦免有三种方式：通过圣洗、通过祈祷、通过更谦卑的补赎；然而天主只赦免受洗者的罪。就连祂起初赦免的罪，也只赦免受洗的人。何时？当他们受洗时。后来因祈祷、因补赎而赦免的罪，祂赦免的对象也是受洗的人。因为尚未成为儿子的人，怎能说\Quote{我们的天父}？慕道者，既为慕道者，就背负他们所有的罪。如果慕道者如此，外邦人岂不更甚？异端者岂不更甚？但我们不更改异端者的圣洗。为什么？因为他们拥有圣洗，就像逃兵拥有兵士的印记一样：这些人也是如此拥有圣洗；他们拥有它，却为此受罚，而非得冠冕。然而，如果逃兵本人改过自新，开始履行兵役，谁敢更改他的印记？\par

17. 我们也信\Quote{肉身的复活}，这已在基督内先行：好让身体也对那在其元首中先行的事怀有希望。基督是教会的元首；教会是基督的身体。我们的元首复活了，升了天：元首在哪里，肢体也在哪里。肉身的复活是怎样的？免得有人偶然以为像拉匝禄的复活那样，为使你知道并非如此，就加上\Quote{进入永生。}愿天主重生于你！愿天主保存并护佑你！愿天主引领你安抵祂自己，祂就是永生。阿们。\par

\SourceRecord{Translated by H. Browne.；Nicene and Post-Nicene Fathers, First Series；Vol. 3.；Edited by Philip Schaff.；Buffalo, NY: Christian Literature Publishing Co.,；1887.；<http://www.newadvent.org/fathers/1307.htm>.}
